\subsection{Settlement Schemes and Vegetation Loss}

Large-scale settlement schemes by the Kenyan government such as the Million Acre Scheme changed land use and land cover (LULC) by converting large-scale farms and forests into smallholder agricultural or subsistence farming plots.
From a macroeconomic perspective, designating an area as a settlement scheme and resettling people from other regions of Kenya \parencite{leoLandClassKenya1984}
such as the reserves implemented by the British colonial government introduces a shock to the local economy and environment.

As mentioned, one of the primary drivers of forest loss and degradation is the LULC change created by converting areas for agricultural production.
\cite{gitauColonialeraSettlementsPostcolonial2025} found that the conversion following Kenya's independence broke up large-scale farms belonging to colonial settlers into smaller plots, thereby reducing the extent of large estates by more than half.
Then, using spatial modeling, they found that the establishment of permanent smallholder settlement resulted in a significant loss of vegetation cover.
Unlike traditional shifting cultivation or pastoralism, permanently allocated plots of land resulted in continuous cultivation, which prevented the natural regeneration of vegetation, according to \cite{serneelsProximateCausesLanduse2001}.
\cite{gitauColonialeraSettlementsPostcolonial2025} attribute roughly 20\% of post-colonial mountain forest loss directly to the agricultural encroachment of smallholder settlement schemes.
The authors also highlight \textquote{a great potential to improve our understanding of long-term trends in tropical forest cover by incorporating historical maps with more recent datasets and methods} \parencite{gitauColonialeraSettlementsPostcolonial2025}, which is the approach taken in this study.

Beyond the main conversion of colonial estates, the establishment on cleared forest area contributed to this loss.
Between 1990 and 2000, following the democratization of Kenya, numerous new settlement schemes were established on degazetted forest land.
\cite{gitauColonialeraSettlementsPostcolonial2025} explicitly stated that, compared to early programs, schemes established in this period saw the conversion of 80,000 ha of primary forest.
Tropical moist forests and montane forests play a critical role as carbon sinks, climate regulators, and biodiversity hotspots by preventing erosion and hosting a variety of floral species \parencite{boeleeEcosystemsWaterFood2011}.
Furthermore, they function as a water tower, capturing rainfall, maintaining a stable groundwater table, and regulating flood and drought cycles \parencite{boeleeEcosystemsWaterFood2011}.
The Mau Forest, where, in terms of total area, most of the settlement schemes were established, feeds 12 major rivers, providing water to over 5 million people \parencite{jacobsLandUseAffects2017}.
Those properties result in very good soil quality, which is why this land was selected \parencite{sanfordDemocratizationElectionsPublic2023}.

\cite{sanfordDemocratizationElectionsPublic2023} found that the Moi government routinely utilized the degazettement and distribution of those schemes as an electoral tool, especially in the Mau Forest region.
The high fertility of the land combined with its availability created immediate patronage profits, without the need for state capital investment.
A report commissioned by the Kenyan government in 2004 documents this systematic and illegal distribution of gazetted forest land \parencite{ndunguReportComissionInquiry2004,southallNdunguReportLand2005}.
\cite{kinyanjuiNDVIbasedVegetationMonitoring2011} showed using the Normalized Difference Vegetation Index (NDVI), that the human encroachment
following the establishment of these settlement schemes caused an abrupt degradation of vegetative health and biodiversity in the Mau Forest region.
However, \cite{kinyanjuiNDVIbasedVegetationMonitoring2011} only analyzed effects for the whole region, as at the time of their study, the dataset of settlement scheme boundaries was not yet available.

As a result of all these factors, demographic changes, land use changes, agricultural clearing, the difference in vegetation cover at the scheme boundaries is expected to be significant,
with the area inside the settlement scheme having a significantly lower vegetation cover than the surrounding area.
Vegetation cover can be measured using the NDVI, which is a widely used remote sensing index used by \cite{khanAssessingImpactLand2025, serneelsProximateCausesLanduse2001,kinyanjuiNDVIbasedVegetationMonitoring2011,lambinMonitoringForestDegradation1999,pettorelliUsingSatellitederivedNDVI2005}.
Additionally, \cite{khanAssessingImpactLand2025} highlights the Normalized Difference Moisture Index (NDMI) to provide an improved understanding of the impact on vegetation moisture content and hydrological conditions,
making it a relevant index in the context of TMF and the Mau Forest region.

Hypothesis 1 (H1): The allocation of land for settlement schemes negatively impacts local environmental conditions, resulting in significantly lower vegetation cover (NDVI and NDMI) within program boundaries compared to neighboring non-program areas.